% ============================================================================
% Section L: Individual Reflections (Deliverable L — Excluded from 12-15 page body)
% Requirement: Each of the 3 members independently writes 200-250 words.
% ============================================================================
\section{Individual Reflections}
\label{sec:L}

\textit{Prompt addressed independently by each group member: ``Which design parameter or later life-cycle consequence was initially outside my view of the assigned system, and how has recognizing it changed the way I would make an engineering decision?''}

\subsection{Reflection 1 --- \memberOneName\ (\texttt{\memberOneRoll})}
Before conducting this practical, I viewed an e-waste management facility primarily through its operational processing flow---specifically how efficiently shredders and optical sorters maximize the Material Recovery Rate (TPM-1). The design parameter entirely outside my initial view was \textbf{DDP-7 (Facility Floor Containment and Design-for-Decommissioning)} and its coupling to the \textbf{Retirement and Disposal Concurrent Life Cycle (LC-4)}. I had implicitly assumed that site decommissioning was a distant administrative formality to be handled in Year~15 once machinery wore out. Analyzing Decision~D3 revealed how a seemingly routine Phase~5 civil cost-saving---pouring standard unsealed concrete with permanently grouted equipment anchors instead of an epoxy-lined secondary containment slab with bolted skid rails---silently locks in irreversible heavy-metal soil seepage and destroys mid-life retrofit flexibility. Saving roughly INR~55 lakh during construction can trigger over INR~2.2~crore in hazardous soil excavation, MPCB closure delays, and zero machinery resale value fifteen years later. Recognizing this consequence has fundamentally altered how I approach engineering design decisions. Rather than optimizing solely for initial capital expenditure or steady-state functional output, I now treat end-of-life disassembly, baseline environmental verification, and decommissioning escrow as active preliminary-design constraints. In any future engineered system I help design, I will explicitly evaluate whether an early physical choice can be cleanly reversed, upgraded, or retired without transferring hidden remediation liabilities to future stakeholders.

\subsection{Reflection 2 --- \memberTwoName\ (\texttt{\memberTwoRoll})}
At the outset of this assignment, my mental model of engineering decision-making focused on maximizing headline technical capability and upfront automation. I initially favored \textbf{Alternative~3 (Monolithic High-Throughput Shred-First Line)} because its 15~t/day throughput (TPM-3) surpassed our 10~t/day criterion by fifty percent while requiring minimal manual labor. What lay outside my view was how external \textbf{Design-Independent Parameters---specifically DIP-2 (Feedstock Product-Mix Shift toward glued Li-ion devices) and DIP-5 (External TSDF Gate Fees)}---interact with internal design choices to invert acquisition-stage economics during utilization. Tracing the full 15-year \textbf{Levelized Life-Cycle Cost (TPM-6)} and running the Year-6 disruption scenario exposed a critical paradox: feeding unsegregated electronics directly into twin shredders cross-contaminates valuable metal and polymer streams, lowers MRR to 79\%, consumes 245~kWh/t of electricity (TPM-5), and amplifies hazardous TSDF disposal volumes whenever regulations tighten. When grid tariffs and TSDF fees rose in Year~6, the supposedly advanced monolithic plant became financially and legally unviable, whereas \textbf{Alternative~2 (Modular Hybrid)} remained resilient through a two-week modular bay upgrade. This insight permanently changed how I evaluate engineering alternatives. I will no longer treat peak nominal performance or single-point cost estimates as decisive; instead, I will stress-test candidate designs against plausible shifts in external DIPs, verify compliance across every mandatory safety and environmental criterion, and prioritize modular life-cycle adaptability over rigid peak throughput.

\subsection{Reflection 3 --- \memberThreeName\ (\texttt{\memberThreeRoll})}
When our group began bounding the UA-EWRRF, I concentrated heavily on regulatory permitting, digital EPR traceability (DDP-5), and initial environmental controls, while largely overlooking \textbf{DDP-3 (Mechanical Maintainability and Critical-Subsystem Redundancy)} and the \textbf{Support and Maintenance Concurrent Life Cycle (LC-3)}. I had treated maintenance as a routine post-commissioning operational task rather than an engineered architectural attribute. Examining how abrasive fiberglass PCBs accelerate shredder-blade wear and how imported optical-sensor lead times govern \textbf{Plant Operational Availability (TPM-2, criterion $A_o \ge 88\%$ with $\text{MTTR} \le 4\text{ h}$)} completely shifted my perspective. Without designing quick-change bolted blade cassettes, overhead gantry clearance, dual parallel baghouse fans, and local spare-parts readiness during Phase~3 and Phase~4, a single jammed rotor or failed dust-collector blower stalls daily throughput (TPM-3), violates CPCB's 180-day storage cap, and pushes work-zone particulates above the $2.0\text{ mg/m}^3$ safety threshold (TPM-4). Furthermore, mapping feedback path \textbf{FB-2} demonstrated that maintenance records are invaluable design intelligence only if contractual channels return failure data to equipment OEMs and future bay engineers. Consequently, my engineering decision-making now incorporates supportability from day one: I will evaluate every subsystem not merely by its rated specification on paper, but by its mean time to repair, local supply-chain resilience, and ability to sustain safe availability across realistic wear cycles.

